\documentclass[a4paper,14pt]{extarticle}
\usepackage{fontspec}
\IfFontExistsTF{Times New Roman}{\setmainfont{Times New Roman}}{\setmainfont{DejaVu Serif}}
\setmonofont{DejaVuSansMono.ttf}
\usepackage[russian]{babel}
\usepackage[left=30mm,right=15mm,top=20mm,bottom=20mm]{geometry}
\usepackage{amsmath,amssymb,booktabs,enumitem,indentfirst,setspace,pgfplots,float,xurl}
\usepackage[labelsep=endash]{caption}
\usepackage[hidelinks]{hyperref}
\hypersetup{pdftitle={Умножение матриц с OpenMP и MPI},pdfauthor={Храмов В. Д.}}
\pgfplotsset{compat=1.18}
\onehalfspacing
\setlength{\parindent}{1.25cm}
\setlength{\emergencystretch}{2em}
\input{build/measurements.tex}

\begin{document}

\begin{titlepage}
\centering
МИНИСТЕРСТВО НАУКИ И ВЫСШЕГО ОБРАЗОВАНИЯ \\
РОССИЙСКОЙ ФЕДЕРАЦИИ\par
\vspace{0.7em}
Санкт-Петербургский политехнический университет Петра Великого\par
Институт компьютерных наук и кибербезопасности\par
Высшая школа технологий искусственного интеллекта\par
Направление 02.03.01 «Математика и компьютерные науки»\par
\vfill
{\Large\bfseries КУРСОВАЯ РАБОТА\par}
\vspace{0.7em}
по дисциплине «Архитектура суперкомпьютеров»\par
\vspace{0.7em}
{\Large\bfseries Умножение матриц с OpenMP и MPI\par}
\vfill
\begin{flushright}
\begin{tabular}{@{}l@{}}
Выполнил студент группы 5130201/40101 \\
Храмов В. Д. \\[1em]
Проверил \\
Чуватов М. В.
\end{tabular}
\end{flushright}
\vfill
Санкт-Петербург\par
2026
\end{titlepage}

\setcounter{page}{2}
\tableofcontents
\clearpage

\section{Постановка задачи}

Цель работы --- исследовать, как распределение вычислений между потоками и процессами влияет на время умножения плотных матриц. По заданию требуется последовательно реализовать четыре режима, сравнить время выполнения и проверить совпадение результатов~\cite{task}.

Для двух квадратных матриц $A,B \in \mathbb{Z}^{N \times N}$ необходимо вычислить
\begin{equation}
  C=AB, \qquad c_{ij}=\sum_{k=0}^{N-1} a_{ik} b_{kj}, \qquad 0\le i,j<N.
  \label{eq:product}
\end{equation}
Все четыре режима используют прямой алгоритм по формуле~\eqref{eq:product}, одни и те же входные данные и одинаковый порядок вычисления каждого элемента.

\begin{table}[H]
\centering
\caption{Конфигурации запуска}
\label{tab:configs}
\small
\begin{tabular}{@{}lccc@{}}
\toprule
Режим & Узлы & Процессы & Потоки на процесс \\
\midrule
Последовательный & 1 & 1 & 1 \\
OpenMP & 1 & 1 & 4 \\
MPI & 2 & 2 & 1 \\
MPI + OpenMP & 2 & 2 & 2 \\
\bottomrule
\end{tabular}
\end{table}

Каждый элемент $C$ требует $N$ умножений и $N$ сложений с аккумулятором, который начинается с нуля. Поэтому программа выполняет $N^3$ скалярных умножений и столько же сложений, а объём данных трёх матриц пропорционален $N^2$. Для вычисления строки $C$ нужны соответствующая строка $A$ и вся матрица $B$. Строки результата независимы, поэтому их удобно распределять между исполнителями.

\clearpage
\section{Реализация}

\subsection{Последовательное ядро и OpenMP}

Матрицы хранятся в одномерных массивах типа \texttt{int32\_t}. Элементы генерируются детерминированно генератором xorshift32 с начальным состоянием 2026 и значениями от $-9$ до $9$. Матрица $B$ сразу записывается в транспонированном виде. При умножении ядро проходит по строке $A$ и строке $B^\mathsf{T}$ с единичным шагом. Элементы обоих массивов читаются последовательно, что улучшает пространственную локальность.

Для настроенного предела $N\le4096$ модуль любого произведения ограничен 81, модуль промежуточной суммы --- $81\cdot4096=331\,776$. Поэтому аккумулятор \texttt{int32\_t} сохраняет точный результат. Исходные матрицы и готовый блок строк представлены одинаково во всех режимах.

OpenMP создаёт команду потоков в области \texttt{parallel}, а директива \texttt{for schedule(static)} распределяет строки $C$ между ними~\cite{openmp}. Потоки читают общие входные массивы и записывают непересекающиеся строки результата. Аккумулятор создаётся внутри тела цикла и используется одним потоком. Статическое распределение подходит для этой задачи, поскольку на каждую строку приходится одинаковый объём работы.

\subsection{MPI и гибридный режим}

В \texttt{MPI\_COMM\_WORLD} каждому процессу присваивается ранг от 0 до $P-1$, где $P$ --- число процессов~\cite{mpi}. Ранг 0 создаёт входные матрицы и выделяет память под итоговую матрицу. При $N=Pq+r$, где $0\le r<P$, первые $r$ рангов получают по $q+1$ строке, остальные --- по $q$. Разность размеров блоков ограничена одной строкой. Размер 257 в эксперименте проверяет случай, когда строки делятся не поровну.

\begin{samepage}
В каждой итерации процессы выполняют следующие действия.
\begin{enumerate}[nosep]
  \item \texttt{MPI\_Bcast} рассылает полную матрицу $B^\mathsf{T}$.
  \item \texttt{MPI\_Scatterv} раздаёт последовательные блоки строк $A$ с учётом числа строк у каждого ранга.
  \item Каждый ранг вычисляет свой блок $C$ общим ядром умножения. В гибридном режиме строки блока делятся между потоками OpenMP.
  \item \texttt{MPI\_Gatherv} собирает блоки $C$ в исходном порядке строк на ранге 0.
\end{enumerate}
\end{samepage}

Гибридный вариант запрашивает уровень \texttt{MPI\_THREAD\_FUNNELED}. Вызовы MPI выполняет поток, инициализировавший MPI, а вычислительная область OpenMP завершает работу перед сбором строк. Четыре исполняемых файла собираются из \path{src/matmul.cpp}. Флаги сборки включают OpenMP и MPI по отдельности или совместно.

При двух рангах передача массива $B^\mathsf{T}$ занимает $4N^2$ байт. Удалённый блок $A$ и его результат $C$ занимают по $4N\lfloor N/2\rfloor$ байт. Для $N=2048$ это 32 МиБ полезных данных на одно умножение. Служебные сообщения и заголовки TCP увеличивают фактический объём трафика.

\section{Методика эксперимента}

\subsection{Стенд}

Эксперимент проведён в Ubuntu 24.04 для ARM64 в Docker Desktop на MacBook Air с Apple M2. Два контейнера \texttt{node-a} и \texttt{node-b} имеют разные сетевые пространства имён и соединены изолированной виртуальной сетью. Для MPI выбран транспорт TCP. Обоим контейнерам доступны одни и те же четыре виртуальных процессора. Последовательный режим и OpenMP выполняются на \texttt{node-a}. Для MPI и гибридного режима запускается по одному рангу на каждом контейнерном узле. Имена узлов и фактическое число потоков записаны в \path{results/placement.txt}.

Сборка выполнена на \texttt{node-a} компилятором GNU C++~13.3.0 с флагами \texttt{-O3} и \texttt{-std=c++17}. Для параллельных вариантов использованы GNU OpenMP и Open MPI~4.1.6~\cite{ompi}. На втором узле запускается копия файла с первого. Паспорт стенда хранится в \path{results/environment.txt}.

\subsection{Измерение времени и правильности}

Для каждого $N\in\{257,512,1024,2048\}$ выполняется один прогрев и семь измерений на конфигурацию. Таймер \texttt{steady\_clock} измеряет время ядра в последовательном режиме и в OpenMP. В MPI-режимах барьер выравнивает начало итерации. После барьера \texttt{MPI\_Wtime} измеряет длительность рассылки, вычислений и сбора на каждом ранге. По окончании итерации \texttt{MPI\_Reduce} выбирает максимум локальных длительностей~\cite{mpi}.

Генерация матриц, выделение основных массивов, старт процессов, проверка и запись CSV вынесены за пределы таймера. Последовательный прогрев формирует эталонную матрицу. Для каждого измерения программа сравнивает все элементы результата с эталоном и при первом несовпадении сообщает координаты. Серии замеров на стенде выполняются по очереди.

Для каждого размера берётся медиана времени $T$ по семи повторам. Ускорение $S$ и эффективность $E$ вычисляются по формулам
\begin{equation}
    S=\frac{T_{\text{посл}}}{T}, \qquad E=\frac{S}{p},
    \label{eq:metrics}
\end{equation}
где $p$ --- число вычислительных исполнителей из таблицы~\ref{tab:configs}. Оно равно 4 для OpenMP и гибридного режима и 2 для MPI. Медианы, минимумы и максимумы серий доступны в \path{results/summary.csv}.

\clearpage
\section{Результаты и обсуждение}

\begin{table}[H]
\centering
\caption{Медианное время умножения, мс}
\label{tab:time}
\small
\begin{tabular}{@{}rrrrr@{}}
\toprule
$N$ & Последовательно & OpenMP & MPI & MPI + OpenMP \\
\midrule
\TimeRows
\bottomrule
\end{tabular}
\end{table}

\begin{figure}[H]
\centering
\begin{tikzpicture}
\begin{axis}[
    width=0.94\textwidth,height=7.1cm,
    xlabel={Размер матрицы $N$},ylabel={Ускорение $S$},
    xtick=data,xticklabel style={/pgf/number format/1000 sep={}},
    ymin=0,ymajorgrids,grid style={gray!25},
    legend style={at={(0.5,-0.23)},anchor=north,legend columns=3,font=\small,draw=none},
    table/col sep=comma,
]
\addplot[blue!70!black,thick,mark=*] table[x=n,y=openmp] {build/speedup.csv};
\addlegendentry{OpenMP}
\addplot[orange!80!black,thick,mark=square*] table[x=n,y=mpi] {build/speedup.csv};
\addlegendentry{MPI}
\addplot[teal!75!black,thick,mark=triangle*] table[x=n,y=hybrid] {build/speedup.csv};
\addlegendentry{MPI + OpenMP}
\end{axis}
\end{tikzpicture}
\caption{Ускорение относительно последовательного режима}
\label{fig:speedup}
\end{figure}

Во всех 112 измерениях результаты совпали с последовательным эталоном. В MPI и гибридном режиме журнал подтвердил по одному процессу на каждом из двух контейнерных узлов. При $N=257$ первый ранг обрабатывал 129 строк, второй --- 128.

% !TEX root = main.tex
\clearpage
\section{Проверка и анализ результатов}

Во всех 112 измерениях получено точное поэлементное совпадение с эталоном. В журнале размещения подтверждено, что MPI-процессы выполнялись на разных контейнерных узлах. При $N=257$ ранг 0 обрабатывал 129 строк, ранг 1~--- 128, что проверяет сбор результата при неравных размерах блоков.

В таблице~\ref{tab:time} приведены медианы семи измерений. Для $N=257$, $512$ и $2048$ наименьшее время получено с OpenMP. При $N=1024$ наименьшая медиана получена в гибридном режиме.

\begin{table}[H]
\centering
\small
\caption{Медианное время выполнения, мс}\label{tab:time}
\begin{tabular}{@{}rrrrr@{}}
\toprule
$N$ & Последовательно & OpenMP & MPI & MPI + OpenMP \\
\midrule
\TimeRows
\bottomrule
\end{tabular}
\end{table}

На рисунке~\ref{fig:speedup} показано ускорение относительно последовательного варианта. Добавление второго потока на ранг увеличило ускорение MPI для всех четырёх размеров при том же объёме передаваемых матриц.

\begin{figure}[H]
\centering
\begin{tikzpicture}
\begin{axis}[
    width=0.94\textwidth,
    height=6cm,
    xlabel={Размер матрицы $N$},
    ylabel={Ускорение $S$},
    xtick=data,
    xticklabel style={/pgf/number format/1000 sep={}},
    ymin=0,
    ymajorgrids,
    grid style={gray!25},
    legend style={at={(0.5,0.06)},anchor=south,legend columns=3,font=\small,draw=none},
    table/col sep=comma,
]
\addplot[blue!70!black,thick,mark=*] table[x=n,y=openmp] {build/speedup.csv};
\addlegendentry{OpenMP}
\addplot[orange!80!black,thick,dashed,mark=square*] table[x=n,y=mpi] {build/speedup.csv};
\addlegendentry{MPI}
\addplot[teal!75!black,thick,dashdotted,mark=triangle*] table[x=n,y=hybrid] {build/speedup.csv};
\addlegendentry{MPI + OpenMP}
\end{axis}
\end{tikzpicture}
\caption{Ускорение относительно последовательного варианта}\label{fig:speedup}
\end{figure}

При $N=1024$ медиана гибридного варианта составила 50,598~мс, а OpenMP~--- 52,964~мс. Диапазоны 50,200--51,712~мс и 49,174--54,998~мс пересекаются. Устойчивое преимущество гибридного варианта на этом размере по одной серии установить нельзя.

При $N=2048$ в режиме MPI получено 867,074~мс, в гибридном~--- 462,048~мс. Сокращение времени составило $46{,}7\,\%$. Для OpenMP с четырьмя потоками получено 444,171~мс, то есть гибридный запуск с тем же общим числом потоков занял примерно на $4{,}0\,\%$ больше времени. По формуле~\eqref{eq:metrics} эффективность составила $94{,}1\,\%$ для OpenMP, $96{,}4\,\%$ для MPI и $90{,}5\,\%$ для гибридного режима. Более высокая эффективность MPI получена при меньшем числе исполнителей.

Время обменов входит в измерения MPI. Разность между временами OpenMP и гибридного варианта включает также влияние планирования потоков и доступа к памяти.


\clearpage
\section{Воспроизведение эксперимента}

Сценарий \path{bench/docker/run.sh} собирает программу в Ubuntu, запускает измерения на двух контейнерных узлах и сохраняет результаты. После работы он удаляет временные контейнеры, сеть и том запуска. Команды выполняются из каталога \path{course_work}.

\begin{verbatim}
bash bench/docker/run.sh
make report
\end{verbatim}

Для запуска на двух Linux-узлах установите совместимые версии Open MPI и настройте доступ по SSH. Соберите варианты командой \texttt{make} на первом узле и скопируйте каталог \texttt{build} на второй по тому же абсолютному пути. С первого узла запустите измерения.

\begin{verbatim}
bash bench/run.sh --host node-a:1,node-b:1
\end{verbatim}

Каталог \path{results} содержит исходные замеры и журнал размещения. Сценарий \path{bench/summarize.py} проверяет их, формирует сводку, таблицу и данные графика для LaTeX.

\renewcommand{\refname}{Список источников}
\begin{thebibliography}{9}
\addcontentsline{toc}{section}{Список источников}
\bibitem{task} Устное задание к курсовой работе по дисциплине «Архитектура суперкомпьютеров» и фотографии доски. Исходные материалы в каталоге \path{course_work/materials}, 2026.
\bibitem{openmp} OpenMP Architecture Review Board. OpenMP Application Programming Interface, версия 5.2. \href{https://www.openmp.org/spec-html/5.2/openmp.html}{Спецификация OpenMP}.
\bibitem{mpi} MPI Forum. MPI: A Message-Passing Interface Standard, версия 4.1. \href{https://www.mpi-forum.org/docs/mpi-4.1/mpi41-report.pdf}{Стандарт MPI}.
\bibitem{ompi} Open MPI Project. \href{https://www.open-mpi.org/doc/v4.1/}{Документация Open MPI 4.1}.
\end{thebibliography}

\end{document}
